% from main.tex
%%%%%%%%%%  scan idx 263  |  folio 250  %%%%%%%%%%
\textbf{Corollaire 3.7.2.} Soit \(f:X\to Y\) un morphisme quasi-compact, quasi-séparé, et localement de tor-dimension finie (3.1), et soit \(E\in\Ob D^b(X)_{\qcoh}\). Si \(E\) est localement d'amplitude plate finie (au sens absolu), il en est de même de \(Rf_*(E)\).

\emph{Preuve.} Conjuguer (3.4.1 a)) et (3.7.1).

\section*{4. Perfection relative}

\textbf{Définition 4.1.} Soient \(f:X\to Y\) un morphisme localement de type fini de schémas, \(E\in\Ob D(X)\), \([a,b]\) un intervalle de \(\mathbb Z\). On dit que \(E\) est d'amplitude parfaite relativement à \(f\) (ou \(Y\) s'il n'y a pas d'ambiguïté), contenue dans \([a,b]\), et l'on écrit
\[
 \parfamp_f(E)\subset[a,b]\qquad(\text{ou }\parfamp_Y(E)\subset[a,b])
\]
si, relativement à \(f\), \(E\) est pseudo-cohérent (1.2) et d'amplitude plate contenue dans \([a,b]\) (3.1). On dit que \(f\) est d'amplitude parfaite contenue dans \([-n,0]\), ou de dimension parfaite \(\leq n\) \((n\geq0)\), si \(\mathcal O_X\) est d'amplitude parfaite contenue dans \([-n,0]\). On dit que \(E\) est d'amplitude parfaite finie relativement à \(f\) s'il existe un intervalle \([a,b]\) tel que l'on ait \(\parfamp_f(E)\subset[a,b]\). On dit que \(E\) est parfait rel. à \(f\) si \(E\) est localement d'amplitude parfaite finie rel. à \(f\). On dit que \(f\) est parfait si \(f\) est localement d'amplitude parfaite finie.

\textbf{Exemple 4.1.1.} Une immersion régulière de codimension \(d\) (VII 1.4) est de dimension parfaite \(\leq d\) (comme on le voit grâce au complexe de Koszul). Un morphisme localement d'intersection complète (i.e. localement composé d'une
